% ============================================================================
% Part 7 --- Faithfulness argument
% ============================================================================
\chapter{The Faithfulness Argument}
\label{ch:faithfulness}

\section{What ``faithful'' means}
\label{sec:faithful-meaning}

\begin{definition}[Faithful lowering]
\label{def:faithful}
The generated Solidity verifier $V_{\mathsf{sol}}$ is a \emph{faithful
translation} of the Rust reference verifier $V_{\mathsf{rust}}$ on profile
$P$ if, for every input $(\mathsf{vk}, \mathsf{proof}, \mathsf{instances})$
in $P$:
\begin{enumerate}[label=(\alph*)]
  \item \textbf{Acceptance agreement:} $V_{\mathsf{rust}}$ accepts iff
        $V_{\mathsf{sol}}$ returns \texttt{true} (otherwise reverts).
  \item \textbf{Semantic checkpoint equality:} at every semantic checkpoint
        (challenge, identity eval, PCS intermediate, pairing input), the
        intermediate value in $V_{\mathsf{sol}}$ equals the corresponding
        value in $V_{\mathsf{rust}}$.
  \item \textbf{Check preservation:} every rejection condition in
        $V_{\mathsf{rust}}$ has a corresponding rejection in
        $V_{\mathsf{sol}}$ that fires on the same inputs.
  \item \textbf{No extra acceptance:} $V_{\mathsf{sol}}$ accepts no input
        that $V_{\mathsf{rust}}$ rejects.
\end{enumerate}
\end{definition}

We cannot prove (a)--(d) for all inputs by testing alone. What we have is a
\emph{layered} argument: a structural correspondence (Chapters
\ref{ch:rust}--\ref{ch:correspondence}) that reduces the claim to a finite
set of semantic checkpoints, plus differential testing that verifies those
checkpoints on a large, adversarial input corpus, plus mutation testing
that verifies the rejection paths.

\section{Semantic checkpoints}
\label{sec:checkpoints}

The lowering is organized around six semantic checkpoints. Each is a point
where the two implementations must agree on a concrete value, and each is
backed by a trace id (or range) in the differential test.

\begin{enumerate}[leftmargin=1.6em,itemsep=3pt]
  \item \textbf{Parser / proof reads.} The proof byte stream is consumed in
        exactly the order of Table~\ref{tab:schedule}; the
        \texttt{proof\_\allowbreak{}cptr == NUM\_\allowbreak{}INSTANCE\_\allowbreak{}CPTR} check pins total
        consumption. Trace ids \boxid{10000+} (commitments) and
        \boxid{20000+} (evals) pin every read value.
  \item \textbf{Transcript challenges.} Every squeeze matches
        Definition~\ref{def:squeeze}. Required ids: \boxid{7--16}.
  \item \textbf{Quotient identities.} Each identity evaluation matches
        \eqref{eq:gate-identity}--\eqref{eq:trash}; the batched numerator
        matches \eqref{eq:nu-y}. Required range \boxid{30000--40000} and
        id \boxid{36}.
  \item \textbf{Linearization.} The expected eval
        $-\nu_y(x)$ and the commitment-side scalars match
        \eqref{eq:linearization}. Id \boxid{34}, selector folds
        \boxid{60000+}.
  \item \textbf{KZG multi-open.} Point sets, $q_{\mathsf{com}}$,
        $f_{\mathsf{eval}}$, $v$, $C_{\mathsf{final}}$ match
        \eqref{eq:qcom}--\eqref{eq:v}. Ranges \boxid{40000--42000},
        ids \boxid{31--33}.
  \item \textbf{Final pairing.} The pairing inputs and result match
        \eqref{eq:final-pairing}. Ids \boxid{27,28,35}.
\end{enumerate}

\section{Differential trace testing}
\label{sec:diff-testing}

The test harness (\cg{test.rs}) runs the Rust verifier with the
\texttt{solidity-verifier-trace} feature and the generated Solidity in a
revm EVM with trace hooks, then calls
\cg{assert\_\allowbreak{}trace\_\allowbreak{}equivalence\_\allowbreak{}and\_\allowbreak{}required\_\allowbreak{}coverage} (L1060):
\begin{lstlisting}[style=rust]
// 1. every Rust trace id present in Solidity
// 2. every Solidity trace id has a Rust oracle
// 3. required coverage: ids 7-16, 31-33, 36, 35,
//    ranges 30000-40000, 40000-41000, 41000-42000, 60000-61000
// 4. byte equality for every id
for (id, solidity_data) in solidity_trace {
    let (name, rust_data) = rust_trace.get(id).unwrap();
    assert_eq!(rust_data, solidity_data,
        "trace mismatch id={id} name={name}");
}
\end{lstlisting}
This is run across the full fixture corpus: poseidon, ivc-keccak,
moonlight-wrap, sha, rsa, hybrid-mt circuits, each with embedded and
separate VK, with and without accumulators.

\section{Mutation and rejection testing}
\label{sec:mutation}

Faithfulness of the \emph{rejection} paths is tested by mutation:
\begin{itemize}[leftmargin=1.4em,itemsep=2pt]
  \item \cg{pbt\_\allowbreak{}solidity\_\allowbreak{}verifies\_\allowbreak{}standard\_\allowbreak{}plonk\_\allowbreak{}embedded\_\allowbreak{}vk\_\allowbreak{}proofs}
        --- positive control.
  \item \cg{pbt\_\allowbreak{}solidity\_\allowbreak{}rejects\_\allowbreak{}wrong\_\allowbreak{}instances} --- flipping an
        instance value must revert (breaks \eqref{eq:instance-eval} and
        the transcript).
  \item \cg{pbt\_\allowbreak{}solidity\_\allowbreak{}rejects\_\allowbreak{}malleated\_\allowbreak{}proofs} --- flipping any
        proof bit must revert (breaks a commitment, an eval, or the
        transcript).
  \item \cg{pbt\_\allowbreak{}solidity\_\allowbreak{}rejects\_\allowbreak{}wrong\_\allowbreak{}verifying\_\allowbreak{}keys} --- a VK
        mismatch must revert (breaks \eqref{eq:vk-digest}).
  \item \cg{pbt\_\allowbreak{}separate\_\allowbreak{}vk\_\allowbreak{}digest\_\allowbreak{}prefix\_\allowbreak{}affects\_\allowbreak{}verification} ---
        the digest prefix binds the circuit.
  \item \cg{pbt\_\allowbreak{}structured\_\allowbreak{}proof\_\allowbreak{}calldata\_\allowbreak{}mutations\_\allowbreak{}are\_\allowbreak{}rejected} ---
        structured (not just bit-flip) mutations revert.
  \item \cg{pbt\_\allowbreak{}proof\_\allowbreak{}layout\_\allowbreak{}repack\_\allowbreak{}and\_\allowbreak{}reader\_\allowbreak{}stay\_\allowbreak{}in\_\allowbreak{}sync} ---
        the repacker and parser agree on layout; truncation reverts.
  \item \cg{pbt\_\allowbreak{}transcript\_\allowbreak{}differential\_\allowbreak{}fuzzer\_\allowbreak{}matches\_\allowbreak{}solidity\_\allowbreak{}trace}
        --- a transcript-level fuzzer matches the Solidity trace.
\end{itemize}

\section{Check-by-check preservation}
\label{sec:check-preservation}

Table~\ref{tab:checks} lists every rejection condition in the Rust
verifier and where the generated Solidity enforces it.

\begin{longtable}{@{}p{0.34\linewidth}p{0.34\linewidth}p{0.22\linewidth}@{}}
\toprule
\textbf{Rust rejection} & \textbf{Solidity enforcement} & \textbf{Status} \\
\midrule
\endfirsthead
\toprule
\textbf{Rust rejection} & \textbf{Solidity enforcement} & \textbf{Status} \\
\midrule
\endhead
\texttt{InvalidInstances} (empty committed) &
  profile-pinned at build time & \chk{pinned} \\
transcript read failure (short proof) &
  \texttt{proof\_\allowbreak{}cptr} exact check + shape guard & \chk{yes} \\
non-canonical scalar ($\geq r$) &
  \texttt{lt(eval, r)} in eval/q\_\allowbreak{}eval loops & \chk{yes} \\
non-canonical G1 (pad bytes / $\geq q$) &
  \sol{common\_\allowbreak{}uncompressed\_\allowbreak{}g1} checks & \chk{yes} \\
identity mismatch ($\nu_y \ne h\cdot Z_H$) &
  pairing fails (linearization) & \chk{yes} \\
PCS opening mismatch &
  final pairing fails & \chk{yes} \\
VK mismatch &
  \texttt{EXPECTED\_\allowbreak{}VK\_\allowbreak{}CODEHASH} + digest & \chk{pinned} \\
wrong instance values &
  transcript divergence $\rightarrow$ pairing fails & \chk{yes} \\
missing precompile &
  constructor smoke test reverts deploy & \chk{yes} \\
\bottomrule
\caption{Every Rust rejection condition has a Solidity counterpart.
``Pinned'' means the condition is excluded by the profile rather than
checked at runtime --- the generated contract cannot be fed a violating
input.}
\label{tab:checks}
\end{longtable}

\subsection{Reverse inventory: every generated revert is justified}
\label{sec:revert-inventory}

Table~\ref{tab:checks} proves \emph{no missing checks} (Rust $\to$
Solidity). The complementary question is whether the generated code
contains any \emph{spurious} or unjustified rejection. The production
\texttt{fee\_12345} verifier has exactly \textbf{34}
\texttt{revert(0,0)} sites; Table~\ref{tab:reverts} groups them and
gives each a Rust justification. Every revert traces to a real
verifier condition or a deployment/VM-integrity guard --- none is
unexplained.

\begin{small}
\begin{longtable}{@{}p{0.26\linewidth}p{3.2em}p{0.30\linewidth}p{0.28\linewidth}@{}}
\toprule
\textbf{Revert group} & \textbf{\#} & \textbf{Enforces} & \textbf{Rust counterpart} \\
\midrule
\endfirsthead
\toprule
\textbf{Revert group} & \textbf{\#} & \textbf{Enforces} & \textbf{Rust counterpart} \\
\midrule
\endhead
Deploy smoke (MCOPY\allowbreak/G1ADD\allowbreak/G1MSM\allowbreak/PAIRING) & 10 & runtime precompiles present \& well-shaped & deploy-time guard (no Rust equiv) \\
ABI-shape head pointers & 1 & dynamic args well-formed & Solidity ABI (no Rust equiv) \\
\texttt{scalar\_inv} zero & 1 & no inverse of $0$ & \texttt{inverse()$\to$None$\to$err} \\
modexp precompile & 2 & \texttt{staticcall}+\texttt{returndatasize} & Rust field inverse (direct) \\
G1 canonicality (pad, $<p$) & 4 & padded/uncompressed canonical & transcript canonicality \\
\texttt{ec\_pairing} result & 1 & pairing equation holds & KZG final pairing \\
VK pin (size+codehash) & 1 & pinned VK runtime & deployment pin (no Rust equiv) \\
VK header cross-check & 1 & \texttt{num\_instances}/\texttt{k}/acc match & codegen consistency \\
ABI envelope (proof/instance len) & 1 & layout matches generated & \texttt{Error::InvalidInstances} \\
\texttt{eval}$<r$ (proof + q evals) & 2 & canonical $\mathbb{F}_r$ & transcript canonicality \\
parser completeness & 1 & \texttt{proof\_cptr} consumed exactly & Rust reads exact layout \\
VM unknown opcode & 1 & \texttt{default} case & VM integrity \\
VM final state & 3 & \texttt{q\_pc==q\_end}, stack empty & VM integrity \\
accumulated \texttt{success} folds & 5 & phase-level failure & (aggregated) \\
\bottomrule
\caption{Reverse inventory: all 34 generated reverts, grouped, each with
a Rust justification. No spurious or unexplained rejection.}
\label{tab:reverts}
\end{longtable}
\end{small}

Together, Tables~\ref{tab:checks} and~\ref{tab:reverts} bracket the
faithfulness claim from both directions: the generated verifier neither
omits a Rust check nor invents an unjustified one.

\section{Fiat--Shamir transcript schedule}
\label{sec:transcript-schedule}

Soundness rests on the transcript: every challenge must be squeezed in
exactly the Rust order over exactly the Rust byte stream, or the
challenges diverge and the pairing fails. Table~\ref{tab:fs-schedule}
is the complete absorb/squeeze schedule, Rust $\leftrightarrow$
generated, with the trace ID each challenge records under.

\begin{small}
\begin{longtable}{@{}clp{0.30\linewidth}p{0.30\linewidth}@{}}
\toprule
\textbf{\#} & \textbf{Op} & \textbf{Rust (\texttt{verifier.rs})} & \textbf{Generated (Yul)} \\
\midrule
\endfirsthead
\toprule
\textbf{\#} & \textbf{Op} & \textbf{Rust (\texttt{verifier.rs})} & \textbf{Generated (Yul)} \\
\midrule
\endhead
1 & absorb & \texttt{common(vk\_\allowbreak{}digest)} & \texttt{common\_word\allowbreak(VK\_DIGEST\_MPTR)} \\
2 & absorb & \texttt{common(instance.len())} & \texttt{common\_word(19)} \\
3 & absorb & \texttt{common(instance[i])} $\times 19$ & \texttt{common\_\allowbreak{}word(inst\_\allowbreak{}be)} $\times 19$ \\
4 & absorb & \texttt{common(advice\_\allowbreak{}com)} (user phase) & \texttt{common\_\allowbreak{}uncompressed\_\allowbreak{}g1} \\
5 & squeeze & user challenge (per phase) & \texttt{squeeze\_\allowbreak{}to(USER\_\allowbreak{}CHAL)} \\
6 & squeeze & $\theta$ \textbf{(id 7)} & \texttt{squeeze\_\allowbreak{}to(THETA\_\allowbreak{}MPTR)} \\
7 & absorb & \texttt{common(lookup\_\allowbreak{}mult\_\allowbreak{}com)} & \texttt{common\_\allowbreak{}uncompressed\_\allowbreak{}g1} \\
8 & squeeze & $\beta$ \textbf{(id 8)} & \texttt{squeeze\_\allowbreak{}to(BETA\_\allowbreak{}MPTR)} \\
9 & squeeze & $\gamma$ \textbf{(id 9)} & \texttt{squeeze\_\allowbreak{}to(GAMMA\_\allowbreak{}MPTR)} \\
10 & absorb & \texttt{common(perm\_com)} & \texttt{common\_\allowbreak{}uncompressed\_\allowbreak{}g1} \\
11 & squeeze & trash \textbf{(id 12)} & \texttt{squeeze\_to\allowbreak(TRASH\_CHALLENGE\_MPTR)} \\
12 & absorb & \texttt{common(trashcan\_\allowbreak{}com)} & \texttt{common\_\allowbreak{}uncompressed\_\allowbreak{}g1} \\
13 & squeeze & $y$ \textbf{(id 10)} & \texttt{squeeze\_\allowbreak{}to(Y\_\allowbreak{}MPTR)} \\
14 & absorb & \texttt{common(quotient\_\allowbreak{}com)} & \texttt{common\_\allowbreak{}uncompressed\_\allowbreak{}g1} \\
15 & squeeze & $x$ \textbf{(id 11)} & \texttt{squeeze\_\allowbreak{}to(X\_\allowbreak{}MPTR)} \\
16 & absorb & \texttt{common(evals)} & \texttt{common\_\allowbreak{}word(eval)} \\
17 & squeeze & $x_1, x_2$ \textbf{(id 13, 14)} & \texttt{squeeze\_\allowbreak{}to(X1\_\allowbreak{}MPTR)}, \texttt{squeeze\_\allowbreak{}to(X2\_\allowbreak{}MPTR)} \\
18 & absorb & \texttt{common(W\_com)} & \texttt{common\_\allowbreak{}uncompressed\_\allowbreak{}g1} \\
19 & squeeze & $x_3$ \textbf{(id 15)} & \texttt{squeeze\_\allowbreak{}to(X3\_\allowbreak{}MPTR)} \\
20 & absorb & \texttt{common(evals)} & \texttt{common\_\allowbreak{}word(eval)} \\
21 & squeeze & $x_4$ \textbf{(id 16)} & \texttt{squeeze\_\allowbreak{}to(X4\_\allowbreak{}MPTR)} \\
\bottomrule
\caption{Complete Fiat--Shamir schedule. The squeeze \emph{order}
matches exactly; the trace \emph{IDs} do not follow it (see note).}
\label{tab:fs-schedule}
\end{longtable}
\end{small}

\paragraph{Audit trap: trace IDs $\neq$ squeeze order.} The challenges
are squeezed in the order $\theta,\beta,\gamma,\text{trash},y,x$ but
record under IDs $7,8,9,\mathbf{12},\mathbf{10},\mathbf{11}$ --- the
trash challenge is squeezed \emph{fourth} yet carries ID 12, while $y$
and $x$ (fifth, sixth) carry 10 and 11. The IDs are a fixed audit
numbering tied to the memory layout, not the transcript order. A
reviewer diffing the two implementations \emph{by trace ID} would see
an apparent reordering; the correct invariant is the \emph{squeeze
sequence} (steps 6--15), which matches exactly. This is worth a
generator-side comment so the numbering is not mistaken for a bug.

\paragraph{Why the schedule is the soundness core.} Every challenge is
a Keccak digest of the transcript prefix (Def.~\ref{def:squeeze}); a
single byte of divergence at any absorb step shifts every subsequent
challenge, and the final pairing (Eq.~\eqref{eq:final-pairing}) fails.
So verifying this 21-step schedule matches --- which the side-by-side
lets a human do in one pass --- is sufficient to conclude the two
implementations share the same Fiat--Shamir structure. The
\texttt{squeeze\_to} reseed rule (digest $\to$ buffer,
\S\ref{sec:transcript-helpers}) is what makes the schedule
self-consistent across both sides.

\section{Deliberate deviations and why they are sound}
\label{sec:deviations}

The lowering is not a line-by-line copy. The deviations are deliberate and
each is justified:

\begin{enumerate}[leftmargin=1.6em,itemsep=3pt]
  \item \textbf{Negated numerator.} The verifier stores $-\nu_y(x)$
        (Definition~\ref{def:neg-numer}) rather than $h(x)$. Sound
        because the $(x^n-1)$ factor is carried on the commitment side as
        $(1-x^n)$; the signs cancel in the pairing. Trace id~36
        re-negates to match Rust's positive numerator.
  \item \textbf{PAIRING LHS/RHS naming.} Historical dual-MSM naming
        (left $=\pi$) is opposite the pairing argument order; the
        generated code passes them swapped, matching
        \eqref{eq:final-pairing-dual}. Documented in the generated
        comment.
  \item \textbf{No curve check at absorption.} The parser checks only the
        EIP-2537 encoding canonicality (Definition~\ref{def:g1-canonical}),
        not on-curve membership. Sound because
        \cg{validate\_\allowbreak{}g1\_\allowbreak{}precompile\_\allowbreak{}coverage} proves every absorbed
        point is later consumed by a G1MSM or pairing, which \emph{do}
        validate membership. A point off-curve makes the pairing fail.
  \item \textbf{Truncated challenges.} $x_3$ (and truncated powers of
        $x_1,x_4$) are masked to 128 bits, mirroring
        \rust{truncated\_\allowbreak{}challenges}. The truncation is part of the
        protocol, not a weakening: the Rust verifier does the same, and
        the differential trace pins the truncated values.
  \item \textbf{Dummy queries.} \texttt{fewer-point-sets} adds eval slots
        to merge point sets, mirroring
        \rust{compute\_\allowbreak{}dummy\_\allowbreak{}queries}. The added evals are transcript
        material and range-checked.
  \item \textbf{Selector buckets.} Simple selectors have no proof eval
        slot; Rust inserts \rust{F::ONE} and the codegen emits a literal
        1 and routes identities to buckets. Algebraically identical to
        \eqref{eq:buckets}.
  \item \textbf{Quotient VM vs.\ straight-line.} The VM is an
        implementation choice for code size; the opcode semantics are
        pinned by the identity-eval trace range \boxid{30000+}.
\end{enumerate}

\section{Residual risks and mitigations}
\label{sec:residual}

\begin{itemize}[leftmargin=1.4em,itemsep=2pt]
  \item \textbf{Profile drift.} If a future circuit violates the profile,
        the generator must refuse. Mitigation: \cg{validate\_\allowbreak{}instance\_\allowbreak{}column\_\allowbreak{}shape}
        and \cg{RotatedInstanceQuery} rejection; the build fails loudly.
  \item \textbf{Trace-id collision.} A new trace id could collide with an
        existing range. Mitigation: the test asserts no unexpected ids.
  \item \textbf{Precompile semantics change.} EIP-2537 edge cases (e.g.\
        identity handling) could differ across forks. Mitigation: pinned
        solc, constructor smoke tests, Prague-semantics documentation.
  \item \textbf{Coverage gaps in fixtures.} A circuit shape not in the
        fixture corpus could hide a bug. Mitigation: the structural
        correspondence (Chapters \ref{ch:rust}--\ref{ch:correspondence})
        is shape-independent; fixtures exercise the lowering, not the
        protocol.
\end{itemize}

\section{Generated-artifact cross-check}
\label{sec:artifact-crosscheck}

Beyond reading the templates, the \emph{actual generated}
\texttt{Halo2Verifier.sol} (poseidon fixture, 150\,KB) was inspected to
confirm the generator emits what the spec claims. Spot-checks:

\begin{itemize}[leftmargin=1.4em,itemsep=1pt]
  \item \textbf{Precompile addresses.} The artifact uses
        \texttt{0x0b} (G1ADD, 6 sites), \texttt{0x0c} (G1MSM, 5 sites),
        \texttt{0x0f} (PAIRING, 2 sites) --- exactly the EIP-2537
        addresses verified in \S\ref{sec:ec-pairing}, confirming the
        addresses are correct (not off-by-one).
  \item \textbf{Pairing layout.} The generated \texttt{ec\_pairing}
        lays out \texttt{[lhs\_\allowbreak{}g1 | G2\_\allowbreak{}BASE | rhs\_\allowbreak{}g1 | NEG\_\allowbreak{}S\_\allowbreak{}G2\_\allowbreak{}BASE]}
        and calls \texttt{staticcall(gas(), 0x0f, scratch, 0x300,
        scratch, 0x20)}, then checks
        \texttt{returndatasize()==0x20 \&\& mload(scratch)}. This is
        $e(\mathsf{lhs},[1]_2)\cdot e(\mathsf{rhs},-[s]_2)=1$ ---
        the \emph{opposite} channel assignment from the Rust
        \rust{DualMSM::check} ($e(\mathsf{left},[s]_2)\cdot
        e(\mathsf{right},-[1]_2)=1$), confirming the documented
        \texttt{PAIRING\_LHS}/\texttt{RHS} naming inversion
        (\S\ref{sec:ec-pairing}) is real in the emitted code, not a
        doc artifact.
  \item \textbf{Truncation mask.} The 128-bit mask
        \texttt{0xffffffffffffffffffffffffffffffff} appears 7 times in
        the artifact (the $x_1$-powers and challenge-truncation sites),
        matching the truncated-challenge finding
        (\S\ref{sec:computations}, \S\ref{sec:parser-audit}).
  \item \textbf{MCOPY staging.} The \texttt{ec\_pairing} uses four
        \texttt{mcopy} calls (Cancun) where the pre-Cancun template used
        mstore chains --- the $\sim$500\,gas/call saving noted in
        \S\ref{sec:ec-pairing} is present in the emitted code.
\end{itemize}

These cross-checks confirm the templates render to the expected
artifact: the spec's claims about the generated code are verified
against the generated code itself, not just the template source.

\section{Summary of the argument}
\label{sec:faith-summary}

\begin{center}
\begin{tikzpicture}[
  node distance=0.5cm and 0.9cm,
  box/.style={draw=black!40, rounded corners=2pt, align=center,
              inner sep=5pt, font=\scriptsize\sffamily},
  arr/.style={-{Stealth[length=2mm]}, black!55, thick}]
\node[box, fill=mathpurple!8] (math) {\textbf{Math spec}\\(Ch.~\ref{ch:math})};
\node[box, fill=rustorange!8, right=of math] (rust) {\textbf{Rust ref}\\implements math};
\node[box, fill=codegenblue!8, right=of rust] (cg) {\textbf{Codegen}\\plans math};
\node[box, fill=soliditygray!8, right=of cg] (sol) {\textbf{Solidity}\\executes plan};
\node[box, fill=tracegreen!10, below=of cg] (trace) {\textbf{Diff trace}\\rust == solidity};
\draw[arr] (math) -- (rust);
\draw[arr] (math) -- (cg);
\draw[arr] (cg) -- (sol);
\draw[arr] (rust) -- (trace);
\draw[arr] (sol) -- (trace);
\end{tikzpicture}
\end{center}

The Rust verifier and the code generator are independently shown to
implement the same mathematics (Chapters \ref{ch:math}--\ref{ch:codegen});
the generated Solidity is shown to execute the plan (Chapter
\ref{ch:solidity}); and the differential trace closes the loop by showing
the two executions agree on every semantic value. Mutation testing covers
the rejection side. Together these constitute a practical, auditable
faithfulness argument for the supported profile.
